\documentclass[10pt]{article}

% ---------------------------------------------------------------
% Kaspar --- one-page capability sheet
% Companion to shadow_pov.tex and fast_send.tex
% ---------------------------------------------------------------

\usepackage[utf8]{inputenc}
\usepackage[T1]{fontenc}
\usepackage[default,semibold]{sourcesanspro}
\usepackage[varqu,varl,scaled=0.94]{inconsolata}
\usepackage{microtype}
\usepackage[a4paper,top=14mm,bottom=13mm,left=17mm,right=17mm]{geometry}
\usepackage{booktabs}
\usepackage{xcolor}
\usepackage{titlesec}
\usepackage{enumitem}
\usepackage{graphicx}
\usepackage{array}
\usepackage{tabularx}
\usepackage[hidelinks]{hyperref}
\usepackage{fancyhdr}

\definecolor{ink}{HTML}{111823}
\definecolor{ink2}{HTML}{46505F}
\definecolor{ink3}{HTML}{6C7787}
\definecolor{rulec}{HTML}{C3CBD6}
\definecolor{accent}{HTML}{1F4E9C}
\definecolor{warm}{HTML}{A82F2B}
\definecolor{band}{HTML}{EEF1F5}

\color{ink}
\renewcommand{\arraystretch}{1.12}
\setlength{\parindent}{0pt}
\setlength{\parskip}{3.2pt}
\setlength{\columnsep}{8mm}

% Section headings: small, letterspaced, with a rule under them.
\titleformat{\section}
  {\normalfont\sffamily\bfseries\footnotesize\color{ink}}
  {}{0pt}{\MakeUppercase}[{\vspace{2pt}\color{rulec}\titlerule[1pt]}]
\titlespacing*{\section}{0pt}{11pt}{5pt}

\pagestyle{fancy}
\fancyhf{}
\renewcommand{\headrulewidth}{0pt}
\fancyfoot[L]{\scriptsize\color{ink3}\sffamily Vincent Mayeski
  \quad\textperiodcentered\quad M2 Tech, Montr\'eal
  \quad\textperiodcentered\quad \href{mailto:mayeski@gmail.com}{mayeski@gmail.com}}
\fancyfoot[R]{\scriptsize\color{ink3}\sffamily
  \href{https://github.com/vincent212/kaspar-hft}{github.com/vincent212/kaspar-hft}
  \quad\textperiodcentered\quad MIT \quad\textperiodcentered\quad \thepage}

% Provenance chip: every performance number says where it came from.
\newcommand{\measured}{{\color{accent}\sffamily\bfseries\scriptsize
  \setlength{\fboxsep}{1.5pt}\fbox{\scshape measured}}}
\newcommand{\code}[1]{{\ttfamily\small #1}}
\newcommand{\tnote}[1]{{\color{ink2}\footnotesize #1\par}}

% One capability tile.
\newcommand{\tile}[2]{%
  {\color{ink3}\sffamily\bfseries\scriptsize\MakeUppercase{#1}}\\[1pt]
  \footnotesize #2\par\vspace{5pt}}

\begin{document}

% ===============================================================
% Masthead
% ===============================================================
{\color{accent}\sffamily\bfseries\scriptsize
 LOW-LATENCY TRADING SYSTEM \quad\textperiodcentered\quad C++20 \quad\textperiodcentered\quad MIT}

\vspace{2pt}
{\sffamily\bfseries\fontsize{40}{40}\selectfont Kaspar-HFT}\\[2pt]
{\color{ink2}\sffamily\bfseries\fontsize{19}{19}\selectfont Open Source}

\vspace{4pt}
{\color{ink2}\sffamily\large A turn-key CME futures trading system and order book simulator.}

\vspace{7pt}
{\color{ink}\hrule height 1.6pt}
\vspace{5pt}

% --- spec rail ---
\noindent\begin{tabularx}{\textwidth}{@{}XXXXX@{}}
{\color{ink3}\sffamily\bfseries\scriptsize VENUE} &
{\color{ink3}\sffamily\bfseries\scriptsize MARKET DATA} &
{\color{ink3}\sffamily\bfseries\scriptsize ORDER ENTRY} &
{\color{ink3}\sffamily\bfseries\scriptsize BOOK} &
{\color{ink3}\sffamily\bfseries\scriptsize LICENCE} \\[-4pt]
\sffamily\bfseries\small CME Globex &
\sffamily\bfseries\small MDP3 SBE &
\sffamily\bfseries\small iLink 3 &
\sffamily\bfseries\small MBP \& MBO L3 &
\sffamily\bfseries\small MIT, open source \\
\end{tabularx}

\vspace{2pt}
{\color{rulec}\hrule height 0.6pt}

\vspace{8pt}
\textbf{From production to simulation.} It is a production trading
system --- multicast market data in, full order books reconstructed order by
order, an execution algorithm on top, order management session out to
the exchange. It is also a an order simulator: the same books rebuilt from recorded
packet captures.

The same strategy code, the same execution algorithm
and the same book run in PCAP replay, in live paper trading, and against the live
exchange. Moving between them is a configuration change, so a backtest
exercises the code path that will trade.

% ===============================================================
\section{What is in the box}

\vspace{-2pt}
{\color{ink3}\footnotesize Feed handler, book, execution logic and exchange session layer all included.}
\vspace{4pt}

\noindent\begin{minipage}[t]{0.485\textwidth}

\tile{Simulation}{Three modes off one binary: \textbf{PCAP replay}, live
multicast paper trading, live iLink execution. Deterministic --- a fixed session
replays to identical fills.}

\tile{Concurrency}{A custom \textbf{C++20 actor framework} --- a formal
concurrency model, and the same code whether the hot path runs across threads or
entirely on one. Scale out by giving actors their own threads; collapse to a
single-threaded hot path by grouping them, from the same source either way.}

\tile{Strategy authoring in C++ or Rust}{C++ in-process actors for the lowest
latency, or Rust across an in-process FFI interop. Either way the strategy runs
in the same process as the book.}

\tile{Execution}{\textbf{Shadow-POV}, a passive/aggressive percentage-of-volume algorithm
that places by following orders in the market.}

\end{minipage}\hfill%
\begin{minipage}[t]{0.485\textwidth}

\tile{Market data}{Full \textbf{MDP3} SBE decoder: incremental book updates,
trades, order-by-order MBO, instrument definitions, snapshot recovery. Sequence
gaps are detected and repaired automatically.}

\tile{Order book}{Order-by-order reconstruction, holding every resting order
individually. That is what makes a queue position, and therefore a fill
inference, definable at all.}

\tile{Order entry}{\textbf{iLink 3} session handler: SBE encoding, HMAC
authentication, sequence management, negotiation and establishment,
primary/secondary failover.}

\end{minipage}

% ===============================================================
\section{Performance}

\vspace{-2pt}
{\color{ink3}\footnotesize AMD EPYC 9374F \quad\textperiodcentered\quad RHEL 9
\quad\textperiodcentered\quad g++ 15 \quad\textperiodcentered\quad
\code{-O3 -march=native}}
\vspace{4pt}

The usual objection to actors is messaging overhead. The answer is
\code{fast\_send}, which runs the receiver's handler inline on the caller's
thread and returns the reply as a value.

\vspace{3pt}
\noindent\begin{tabularx}{\textwidth}{@{}Xr>{\raggedright\arraybackslash}p{62mm}@{}}
\toprule
{\color{ink3}\sffamily\bfseries\scriptsize PATH} &
{\color{ink3}\sffamily\bfseries\scriptsize P50} &
{\color{ink3}\sffamily\bfseries\scriptsize WHAT IT COSTS} \\
\midrule
\code{fast\_send} --- inline, same thread & {\color{accent}\bfseries\code{\textasciitilde10 ns}} & $\approx$ one polymorphic virtual call \\
Queued send, one \code{Group} thread & \code{\textasciitilde90 ns} & queue push/pop + dispatch \\
Queued send, separate threads & \code{\textasciitilde3{,}370 ns} & cross-core mailbox wakeup, twice \\
\midrule
Pooled round trip (2 messages) & \code{\textasciitilde92 ns} & vs \code{\textasciitilde128 ns} with plain \code{new} \\
Allocator tail (max observed) & \code{\textasciitilde33 \textmu s} & vs \code{\textasciitilde5.5 ms} unpooled --- the durable win \\
\bottomrule
\end{tabularx}

\vspace{3pt}
\tnote{\measured\ Round trip through the framework --- send, handle, reply ---
one message in flight, from the microbenchmarks in \code{actors/cpp/perf}.}


% ===============================================================
\end{document}
